% !TEX root = chapter-standalone.tex


\section[Linear algebra]{Linear algebra}
\label{sec:cats-of-linear-algebra}

\todotext{J: @J: add more material to this section}
\todotext{@JL: Shouldn't we have just $\natnumbers$ as the objects?}
\begin{ctdefinition}[Category of real matrices]
    \SYNDEF{category of real matrices}
    \label{def:cat-of-real-matrices}
    The category \MatR of real matrices is:
    \begin{enumerate}
        % \item \emph{Objects}: $\styleobj{\reals^n}$, for $n \setin \natnumbers$.
        \item \emph{Objects}: natural numbers $\natnumbers$.
              % \item \emph{Morphisms}: for any $n, m \setin \natnumbers$, $\HomSet{\Set}{\styleobj{\reals^m}}{\styleobj{\reals^n}}$ is the set  of $n \times m$ real matrices.
        \item \emph{Morphisms}: for any $m, n \setin \natnumbers$, $\HomSet{\MatR}{m}{n}$ is the set  of $n \times m$ real matrices.
        \item \emph{Composition}: \SY{matrix multiplication}.
        \item \emph{Identity morphisms}: \SY{identity matrices}.
    \end{enumerate}
\end{ctdefinition}

\begin{ctdefinition}[Category of vector spaces over a field]
    \SYNDEF{category of vector spaces}
    \label{def:cat-of-vector-spaces}
    \label{def:cat-of-real-vector-spaces}
    Let~$\field$ be a field.
    The category~$\VectCat{\field}$ of vector spaces over~$\field$ is:
    \begin{enumerate}
        \item \emph{Objects}: all vector spaces over~$\field$.
        \item \emph{Morphisms}: $\HomSet{\VectCat{\field}}{\Obja}{\Objb}$ is the set of $\field$-linear maps $\Obja \mto \Objb$.
        \item \emph{Composition}: the usual composition of linear maps.
        \item \emph{Identity morphisms}: for any vector space~$\Obja = \tup{\stylesets{X}, +,{{ \cdot}}}$ over~$\field$, the identity morphism $\catid_\Obja$ is given by the identity function $\stylesets{X} \to \stylesets{X}$.
    \end{enumerate}
\end{ctdefinition}

\begin{ctdefinition}[Category of vector spaces and linear relations]
    \SYNDEF{category of linear relations}
    \label{def:cat-of-linear-relations}
    Let~$\field$ be a field.
    The category~$\stylecat{LinRel}_\field$ of vector spaces over~$\field$ and linear relations is:
    \begin{enumerate}
        \item \emph{Objects}: all vector spaces over~$\field$.
        \item \emph{Morphisms}: a morphism from~$\Obja$ to~$\Objb$ is a \emph{linear relation}, i.e., a linear subspace~$\relA \subseteq \Obja \oplus \Objb$.
        \item \emph{Composition}: given linear relations~$\relA \subseteq \Obja \oplus \Objb$ and~$\relB \subseteq \Objb \oplus \Objc$, their composite is
              \begin{equation}
                  \relA \mthen \relB = \makeset{ \tup{\ela, \elc} \setin \Obja \oplus \Objc \mid \exists\, \elb \setin \Objb \colon \tup{\ela, \elb} \setin \relA \text{ and } \tup{\elb, \elc} \setin \relB }.
              \end{equation}
        \item \emph{Identity morphisms}: for any vector space~$\Obja$, the identity morphism is the diagonal subspace~$\makeset{ \tup{\ela, \ela} \mid \ela \setin \Obja } \subseteq \Obja \oplus \Obja$.
    \end{enumerate}
\end{ctdefinition}

\begin{ctdefinition}[Category of Euclidean spaces and isometries]
    \SYNDEF{category of Euclidean spaces and isometries}
    \label{def:cat-of-euclidean-isometries}
    A \emph{Euclidean space} is a finite-dimensional real vector space~$\Obja$ equipped with an inner product, i.e., a bilinear form~$\tup{\,\cdot\,,\,\cdot\,}_\Obja \colon \Obja \times \Obja \to \reals$ that is symmetric, positive definite, and nondegenerate.

    An \emph{isometry} between Euclidean spaces~$\tup{\Obja, \tup{\,\cdot\,,\,\cdot\,}_\Obja}$ and~$\tup{\Objb, \tup{\,\cdot\,,\,\cdot\,}_\Objb}$ is a linear map~$\mora \colon \Obja \mto \Objb$ that preserves the inner product:
    \begin{equation}
        \tup{\mora(\ela),\, \mora(\elb)}_\Objb = \tup{\ela,\, \elb}_\Obja \quad \text{for all } \ela, \elb \setin \Obja.
    \end{equation}

    The category~$\stylecat{Euc}$ of Euclidean spaces and isometries is:
    \begin{enumerate}
        \item \emph{Objects}: all Euclidean spaces.
        \item \emph{Morphisms}: isometries.
        \item \emph{Composition}: composition of linear maps.
        \item \emph{Identity morphisms}: identity functions.
    \end{enumerate}
\end{ctdefinition}



\begin{ctdefinition}[Category of orthogonal spaces and isometries]
    \SYNDEF{category of orthogonal spaces and isometries}
    \label{def:cat-of-orthogonal-isometries}
    Let~$\field$ be a field.
    An \emph{orthogonal space} over~$\field$ is a finite-dimensional vector space~$\Obja$ over~$\field$ equipped with a symmetric bilinear form~$\tup{\,\cdot\,,\,\cdot\,}_\Obja \colon \Obja \times \Obja \to \field$ that is nondegenerate.

    An \emph{isometry} between orthogonal spaces~$\tup{\Obja, \tup{\,\cdot\,,\,\cdot\,}_\Obja}$ and~$\tup{\Objb, \tup{\,\cdot\,,\,\cdot\,}_\Objb}$ is a linear map~$\mora \colon \Obja \mto \Objb$ satisfying
    \begin{equation}
        \tup{\mora(\ela),\, \mora(\elb)}_\Objb = \tup{\ela,\, \elb}_\Obja \quad \text{for all } \ela, \elb \setin \Obja.
    \end{equation}

    The category~$\stylecat{Orth}_\field$ of orthogonal spaces and isometries over~$\field$ is:
    \begin{enumerate}
        \item \emph{Objects}: all orthogonal spaces over~$\field$.
        \item \emph{Morphisms}: isometries.
        \item \emph{Composition}: composition of linear maps.
        \item \emph{Identity morphisms}: identity functions.
    \end{enumerate}
\end{ctdefinition}

\begin{remark}
    A Euclidean space is an orthogonal space over~$\reals$ whose symmetric bilinear form is additionally positive definite.
    Hence the category~$\stylecat{Euc}$ is a \SY{subcategory} of~$\stylecat{Orth}_\reals$.
\end{remark}

\devel{

\begin{ctdefinition}[Category of finite-dimensional Hilbert spaces and isometries]
    \SYNDEF{category of Hilbert spaces and isometries}
    \label{def:cat-of-hilbert-isometries}
    A \emph{finite-dimensional Hilbert space} is a finite-dimensional complex vector space~$\Obja$ equipped with an inner product, i.e., a map~$\tup{\,\cdot\,,\,\cdot\,}_\Obja \colon \Obja \times \Obja \to \mathbb{C}$ that is linear in the second argument, conjugate-linear in the first argument (i.e., \emph{sesquilinear}), Hermitian (i.e.,~$\tup{\ela, \elb}_\Obja = \overline{\tup{\elb, \ela}_\Obja}$), and positive definite (i.e.,~$\tup{\ela, \ela}_\Obja > 0$ for all~$\ela \neq 0$).

    An \emph{isometry} between finite-dimensional Hilbert spaces~$\tup{\Obja, \tup{\,\cdot\,,\,\cdot\,}_\Obja}$ and~$\tup{\Objb, \tup{\,\cdot\,,\,\cdot\,}_\Objb}$ is a linear map~$\mora \colon \Obja \mto \Objb$ that preserves the inner product:
    \begin{equation}
        \tup{\mora(\ela),\, \mora(\elb)}_\Objb = \tup{\ela,\, \elb}_\Obja \quad \text{for all } \ela, \elb \setin \Obja.
    \end{equation}

    The category~$\stylecat{Hilb}_{\mathrm{fd}}$ of finite-dimensional Hilbert spaces and isometries is:
    \begin{enumerate}
        \item \emph{Objects}: all finite-dimensional Hilbert spaces.
        \item \emph{Morphisms}: isometries.
        \item \emph{Composition}: composition of linear maps.
        \item \emph{Identity morphisms}: identity functions.
    \end{enumerate}
\end{ctdefinition}

\begin{remark}
    A Euclidean space can be seen as a finite-dimensional Hilbert space whose underlying vector space is real (rather than complex), with a bilinear (rather than sesquilinear) inner product.
\end{remark}

\todotext{Not sure if it makes sense to include the two symplectic definitions below... only needed if we want to talk about applications involving them...}

\begin{ctdefinition}[Category of symplectic vector spaces and isometries]
    \SYNDEF{category of symplectic spaces and isometries}
    \label{def:cat-of-symplectic-isometries}
    Let~$\field$ be a field.
    A \emph{symplectic vector space} over~$\field$ is a finite-dimensional vector space~$\Obja$ over~$\field$ equipped with a bilinear form~$\omega_\Obja \colon \Obja \times \Obja \to \field$ that is \emph{skew-symmetric} (i.e.,~$\omega_\Obja(\ela, \elb) = -\omega_\Obja(\elb, \ela)$ for all~$\ela, \elb$) and \emph{nondegenerate}.

    A \emph{symplectic isometry} (also called a \emph{symplectomorphism}) between symplectic vector spaces~$\tup{\Obja, \omega_\Obja}$ and~$\tup{\Objb, \omega_\Objb}$ is a linear map~$\mora \colon \Obja \mto \Objb$ satisfying
    \begin{equation}
        \omega_\Objb(\mora(\ela),\, \mora(\elb)) = \omega_\Obja(\ela,\, \elb) \quad \text{for all } \ela, \elb \setin \Obja.
    \end{equation}

    The category~$\stylecat{Symp}_\field$ of symplectic vector spaces and isometries over~$\field$ is:
    \begin{enumerate}
        \item \emph{Objects}: all symplectic vector spaces over~$\field$.
        \item \emph{Morphisms}: symplectic isometries.
        \item \emph{Composition}: composition of linear maps.
        \item \emph{Identity morphisms}: identity functions.
    \end{enumerate}
\end{ctdefinition}



\begin{ctdefinition}[Category of symplectic vector spaces and Lagrangian relations]
    \SYNDEF{category of Lagrangian relations}
    \label{def:cat-of-lagrangian-relations}
    Let~$\field$ be a field.
    Given a symplectic vector space~$\tup{\Obja, \omega_\Obja}$ over~$\field$, a subspace~$\stylesets{L} \subseteq \Obja$ is called \emph{Lagrangian} if~$\stylesets{L} = \stylesets{L}^\perp$, where
    \begin{equation}
        \stylesets{L}^\perp = \makeset{ \ela \setin \Obja \mid \omega_\Obja(\ela, \elb) = 0 \text{ for all } \elb \setin \stylesets{L} }.
    \end{equation}

    Given two symplectic vector spaces~$\tup{\Obja, \omega_\Obja}$ and~$\tup{\Objb, \omega_\Objb}$, a \emph{Lagrangian relation} from~$\Obja$ to~$\Objb$ is a Lagrangian subspace of the symplectic vector space
    $$\tup{\Obja \oplus \Objb,\, -\omega_\Obja \oplus \omega_\Objb}.$$

    The category~$\stylecat{LagRel}_\field$ of symplectic vector spaces and Lagrangian relations over~$\field$ is:
    \begin{enumerate}
        \item \emph{Objects}: all symplectic vector spaces over~$\field$.
        \item \emph{Morphisms}: a morphism from~$\Obja$ to~$\Objb$ is a Lagrangian relation from~$\Obja$ to~$\Objb$.
        \item \emph{Composition}: given Lagrangian relations~$\relA \subseteq \Obja \oplus \Objb$ and~$\relB \subseteq \Objb \oplus \Objc$, their composite is
              \begin{equation}
                  \relA \mthen \relB = \makeset{ \tup{\ela, \elc} \setin \Obja \oplus \Objc \mid \exists\, \elb \setin \Objb \colon \tup{\ela, \elb} \setin \relA \text{ and } \tup{\elb, \elc} \setin \relB }.
              \end{equation}
        \item \emph{Identity morphisms}: for any symplectic vector space~$\Obja$, the identity morphism is the diagonal~$\makeset{ \tup{\ela, \ela} \mid \ela \setin \Obja }$, which is a Lagrangian subspace of~$\tup{\Obja \oplus \Obja,\, -\omega_\Obja \oplus \omega_\Obja}$.
    \end{enumerate}
\end{ctdefinition}

\todotext{@JL: Verify that the composition of Lagrangian relations is again Lagrangian (this is a nontrivial fact). Add a remark or reference.}

}

